\documentclass[]{article}

\usepackage{amsmath}
\usepackage{multirow}
\usepackage{natbib}
\usepackage{graphicx} 


\begin{document}

\subsection{Data and Feature Construction}
The primary dataset consists of daily closing prices for nine major US sector ETFs (XLK, XLF, XLV, XLE, XLI, XLY, XLP, XLU, XLB) and the CBOE Volatility Index (VIX) from January 2015 to April 2026. The S\&P 500 index (\textasciicircum GSPC) was also included as a market-wide control variable. All data was processed to ensure a balanced panel by removing days with missing observations for any of the assets.

The core metric, Sectoral Volatility Dispersion (SVD), was constructed through a multi-step process. First, daily logarithmic returns were calculated for each of the nine sector ETFs. Second, for each ETF, we computed a 21-day rolling realized volatility, defined as the standard deviation of its log-returns over the preceding 21 trading days. This daily volatility figure was then annualized:
\begin{equation}
    \sigma_{i,t}^{(21)} = \sqrt{252} \times \text{StDev}(r_{i, t-20:t})
\end{equation}
where $\sigma_{i,t}^{(21)}$ is the annualized 21-day realized volatility for sector $i$ at day $t$, and $r_{i, t-20:t}$ is the series of daily log-returns for that sector over the 21-day window.

The SVD for each day $t$ was then calculated as the cross-sectional standard deviation of these nine concurrent realized volatility estimates:
\begin{equation}
    \text{SVD}_t = \sqrt{\frac{1}{N-1} \sum_{i=1}^{N} (\sigma_{i,t}^{(21)} - \bar{\sigma}_{t}^{(21)})^2}
\end{equation}
where $N=9$ is the number of sectors and $\bar{\sigma}_{t}^{(21)}$ is the cross-sectional average of the sector volatilities at time $t$. To ensure stationarity and facilitate interpretation in our regression models, this raw SVD series was transformed into a normalized Z-score using a rolling 252-day (one trading year) lookback window for the mean and standard deviation. Stationarity of the resulting time series was confirmed using the Augmented Dickey-Fuller (ADF) test.

To capture the market's correlational structure, we also computed a daily measure of average cross-sector correlation. This was calculated as the mean of all unique pairwise correlations across the nine sector ETFs, based on their daily returns over a rolling 21-day window.

\subsection{VIX Regime Identification}
To endogenously classify market volatility states, we employed a three-state Gaussian Hidden Markov Model (HMM). This probabilistic approach avoids the use of arbitrary, fixed thresholds for the VIX and instead allows the data to define distinct regimes based on the distributional properties of VIX changes. The model was fitted to the daily log-returns of the VIX index, $\Delta \ln(\text{VIX}_t) = \ln(\text{VIX}_t / \text{VIX}_{t-1})$, to ensure the input series was stationary. The three latent states identified by the HMM were subsequently labeled as ``Low'', ``Moderate'', and ``High'' volatility regimes by examining the average raw VIX level corresponding to the days classified within each state.

\subsection{Econometric Analysis}
Our empirical investigation involved two primary modeling frameworks to assess the predictive power of SVD.

First, to test the conditional relationship between SVD and future VIX movements, we estimated an Ordinary Least Squares (OLS) regression model. The model predicts the $k$-day ahead log-return of the VIX using the normalized SVD, an interaction term between SVD and the average sector correlation, and several control variables. The full specification is:
\begin{equation}
    \Delta \ln(\text{VIX})_{t+k} = \alpha + \beta_1 \text{SVD}_t + \beta_2 (\text{SVD}_t \times \text{AvgCorr}_t) + \beta_3 \Delta \ln(\text{VIX})_t + \beta_4 r_{\text{MKT},t} + \epsilon_t
\end{equation}
where $\Delta \ln(\text{VIX})_{t+k}$ is the log-return of the VIX over the next $k$ days, $\text{SVD}_t$ is the normalized SVD at time $t$, $\text{AvgCorr}_t$ is the 21-day average sector correlation, $\Delta \ln(\text{VIX})_t$ is the concurrent VIX log-return, and $r_{\text{MKT},t}$ is the concurrent S\&P 500 log-return. We estimated this model for horizons of $k=5$ and $k=21$ days. To ensure robust statistical inference in the presence of heteroskedasticity and autocorrelation inherent in overlapping time-series data, all reported standard errors were corrected using the Newey-West estimator.

Second, to evaluate SVD as a discrete warning signal for regime shifts, we used a logistic regression model. The model was designed to predict the probability of the market transitioning into the ``High'' volatility regime on day $t$, conditional on not being in that state on day $t-1$. The primary predictor was the lagged normalized SVD ($SVD_{t-1}$). Model performance was evaluated using two metrics: the Area Under the Receiver Operating Characteristic curve (AUROC), which measures overall discriminatory power, and the Precision-Recall Area Under the Curve (PR AUC), which is more informative for evaluating classifiers on highly imbalanced datasets where the event of interest (a transition to high volatility) is rare.

Finally, to confirm that our findings were not driven by the behavior of any single sector, we performed a Leave-One-Out Cross-Validation (LOOCV) on the logistic regression model. The SVD metric was recalculated nine times, each time excluding one sector ETF, and the AUROC was computed for each iteration to assess the stability of the model's predictive performance.

\end{document}
                